\documentclass[11pt]{article}
\usepackage[a4paper,margin=18mm]{geometry}
\usepackage{fontspec}
\setmainfont{Latin Modern Roman}
\usepackage{amsmath,amssymb,amsthm,amscd,graphicx,color}
\usepackage{xurl}
\usepackage[unicode,hidelinks]{hyperref}
\allowdisplaybreaks
\setlength\emergencystretch{3em}
\numberwithin{equation}{section}

\newtheorem{Theorem}{Theorem}[section]
\newtheorem{Corollary}[Theorem]{Corollary}
\newtheorem{Lemma}[Theorem]{Lemma}
 { \theoremstyle{definition}
\newtheorem{Definition}[Theorem]{Definition}
\newtheorem{Remark}[Theorem]{Remark} }

\def\ta{\theta}
\def\Ta{\Theta}
\newcommand{\ti}{\mathrm i}
\newcommand\bn{\mathbf n}
\newcommand\bk{\mathbf k}
\newcommand\bl{\mathbf l}
\newcommand\bx{\mathbf x}

\makeatletter\newlabel{amit}{{3.1}{}{Original source locator}{}{}}\makeatother
\makeatletter\newlabel{bcmit}{{3.4}{}{Original source locator}{}{}}\makeatother
\begin{document}
\begin{center}\Large For rank r>=2 and a nonzero elliptic nome p of absolute value less than one, the coefficient-regular lower-triangular C\_r theta-product matrices of Theorem 3.5 are mutually inverse\end{center}
\noindent\textbf{Stable ID:} 1621\par
\noindent\textbf{Authors:} Hjalmar Rosengren; Michael J. Schlosser\par
\noindent\textbf{Original work:} Multidimensional Matrix Inversions and Elliptic Hypergeometric Series on Root Systems\par
\noindent\textbf{Version:} Official SIGMA 16 (2020) article 088, DOI 10.3842/SIGMA.2020.088; journal PDF and native TeX acquired 2026-09-30, exact bytes pinned by SHA256\par
\noindent\textbf{Selected task scope:} Exactly the elliptic Theorem3.5 for r>=2, 0<|p|<1, nonzero b and nonzero complex parameter sequences c\_j(t), 1<=j<=r, t in Z, at coefficient-regular choices where every theta denominator in the finite relations under consideration is nonzero. The lower-triangular multi-index matrices f\_nk and g\_kl are precisely the printed theta products; inversion means sum\_\{l<=k<=n\} f\_nk g\_kl=delta\_nl (and the equivalent reverse relation) with coordinatewise integer ordering and finite sums. No singular-parameter continuation, r=1 or p=0 degeneration, or separate series-summation theorem is claimed.\par
\noindent\textbf{Difficulty:} H2: The new obligation is a multidimensional inverse relation for arbitrary parameter sequences, whose product coupling prevents the simpler theta-zero-counting argument used for the neighboring inversions. The proof converts a finite rectangular sum through a carefully chosen theta partial-fraction identity and supplies an explicit antisymmetric pairing that cancels the residual sums. This is substantially beyond a one-variable interpolation exercise or a direct application of a known inverse matrix.\par
\noindent\textbf{Editorial boundary note (not author text):} The full original Theorem 3.5 is retained verbatim; the stable item's selected task is its coefficient-regular elliptic rank-at-least-two branch. The complete new author cancellation proof through Lemma 3.6 is included with theta notation, finite inverse-matrix definition, Gustafson's explicitly cited standard identity, the precise parameter specialization and the residual antisymmetric pairing. Neighboring inversion results and summation applications are uncounted and excluded. The existing author gamma-index slip is retained verbatim and flagged separately.\par
\noindent\textbf{Source:} \url{https://sigma-journal.com/2020/088/}\quad DOI: \url{https://doi.org/10.3842/SIGMA.2020.088}\par
\noindent\textbf{License and attribution:} Copyright retained by the named authors. Published by SIGMA under CC BY 4.0: \url{https://creativecommons.org/licenses/by/4.0/}. Source: \url{https://sigma-journal.com/about.html}. Adaptation consists of standalone excerpt formatting, cover and numbering restoration; original author language and mathematical text are retained.\par
\noindent\textbf{Editorial symbol notice (not author text):} Editorial symbol notice, separate from the author: the final antisymmetric expression in Lemma3.6 (PDF p.10) defines gamma using c\_r(k\_r), although k\_r has already been summed out in (3.2). The preceding explicitly printed A\_\{k\_1,...,k\_(r-1)\} and C\_l formulas involve only the outer indices and specify the b/(c\_1(l)c\_1(k\_1)...c\_(r-1)(k\_(r-1))) product, fixing gamma as the corresponding product independent of k\_1,l. The unchanged source retains the index slip; no new cancellation proof is inserted. Parent should retain this warning with the selected original pages.\par
\noindent\textbf{External original locator amit:} Original Theorem 3.1 is the neighboring elliptic A\_r inversion, mentioned only for comparison.\par
\noindent\textbf{External original locator bcmit:} Original Theorem 3.4 is the neighboring elliptic BC\_r inversion, mentioned only for comparison.\par
\medskip\hrule\medskip\noindent\textbf{Author text begins below.}\par
\setcounter{section}{1}
\setcounter{subsection}{0}
\setcounter{subsubsection}{0}
\setcounter{Theorem}{0}
\setcounter{equation}{1}
% Original source lines 116--159
\section{Preliminaries}

\subsection{Theta functions}

We recall some classical results for theta functions.
We will write
\begin{gather*}
\ta(x;p)=\prod\limits_{j=0}^\infty\big(1-p^jx\big)\big(1-p^{j+1}/x\big),
\end{gather*}
where $p$ is a complex number with $|p|<1$.
We refer to the case
 $p=0$, $\ta(x;0)=1-x$, as the \emph{trigonometric case}.
We will also use the shorthand notation
\begin{gather*}
\theta(x_1,\dots,x_n;p)=\ta(x_1;p)\dotsm\ta(x_n;p),
\end{gather*}
\begin{gather*}
\theta\big(x_1y^{\pm},\dots,x_ny^{\pm};p\big)=\ta(x_1y;p)\ta(x_1/y;p)\dotsm\ta(x_ny;p)
\ta(x_n/y;p).
\end{gather*}
In Section~\ref{gmis}, we will suppress the dependence on $p$ and simply write
$\theta(x)=\theta(x;p)$.


Among the properties of $\theta$ we mention the inversion formula
\begin{gather}
\label{ti}\theta(1/x;p)=-\frac1 x \theta(x;p),
\end{gather}
the quasi-periodicity
\begin{gather}
\label{qpt}\ta(px;p)=-\frac1 x \theta(x;p)
\end{gather}
and Weierstrass' addition formula
\begin{gather}\label{tadd}
\ta(xy,x/y,uv,u/v;p)-\ta(xv,x/v,uy,u/y;p)=\frac uy \ta(yv,y/v,xu,x/u;p).
\end{gather}
We will also need the identity
\begin{gather}\label{cpf}
\sum_{l=1}^k\frac{a_l\prod\limits_{j=1}^{k-2}\theta\big(a_lb_j^\pm;p\big)}
{\prod\limits_{j=1,\,j\neq l}^k\theta\big(a_la_j^\pm;p\big)}=0, \qquad
k\geq 2.
\end{gather}
To our knowledge, it was first obtained by Gustafson
\cite[Lemma~4.14]{gu}, see~\cite{r} for further references and comments.

\setcounter{section}{2}
\setcounter{subsection}{2}
\setcounter{subsubsection}{0}
\setcounter{Theorem}{2}
\setcounter{equation}{6}
% Original source lines 271--290
\section{Elliptic multidimensional matrix inversions}\label{mis}


\subsection{General form of the matrix inversions}
\label{gmis}


We will consider matrices
 $f=(f_{\bf nk})_{{\bf n}, {\bf k}\in\mathbb Z^r}$
labelled by pairs of multi-indices. All matrices that appear are
lower-triangular in the sense that $f_{\bf nk}=0$ unless $\bf
n\geq\bf k$, that is, $n_i\geq k_i$ for all $i$. Then,
$f$ and $g=(g_{\bf nk})_{{\bf n}, {\bf k}\in\mathbb Z^r}$
are inverse matrices if
\begin{gather}\label{fg}
\sum_{\bf l\leq k\leq n}f_{\bf nk}g_{\bf kl}=\delta_{\bf nl} \end{gather}
or, equivalently,
\begin{gather*}
\sum_{\bf l\leq k\leq n}g_{\bf nk}f_{\bf kl}=\delta_{\bf nl}.
\end{gather*}

\setcounter{section}{3}
\setcounter{subsection}{1}
\setcounter{subsubsection}{0}
\setcounter{Theorem}{0}
\setcounter{equation}{1}
% Original source lines 294--295
 elliptic $A_r$ matrix inversion. Here, and in the remainder of Section~\ref{gmis}, we write
 $\theta(x)=\theta(x;p)$, where $p$ is viewed as fixed.

\setcounter{section}{3}
\setcounter{subsection}{1}
\setcounter{subsubsection}{0}
\setcounter{Theorem}{4}
\setcounter{equation}{1}
% Original source lines 464--643
Finally, we give a matrix inversion that we associate with the
root system $C_r$. The case $p=0$ of Theorem~\ref{cmit} is \cite[Theorem~4.1]{s}.

\begin{Theorem}[an elliptic $C_r$ matrix inversion]\label{cmit}
Let $(c_j(t))_{t\in\mathbb Z}$, $1 \leq j \leq r$, and $b$ be arbitrary scalars. Then the lower-triangular matrices
\begin{gather*}f_{\bf nk}=\prod\limits_{i=1}^r\frac{\prod\limits_{t=k_i}^{n_i-1}
\bigg\{\theta(c_i(t)b/c_1(k_1)\dotsm c_r(k_r))\prod\limits_{j=1}^r\theta(c_i(t)c_j(k_j))\bigg\}}
{\prod\limits_{t=k_i+1}^{n_i} \bigg\{\theta(c_i(t)c_1(k_1)\dotsm c_r(k_r)/b)
\prod\limits_{j=1}^r\theta(c_i(t)/c_j(k_j))\bigg\}}
\end{gather*}
and
\begin{gather*}
g_{\bf kl}=
\prod\limits_{1\leq i<j\leq r}\frac{\theta\big(c_j(l_j)c_i(l_i)^\pm\big)}
{\theta\big(c_j(k_j)c_i(k_i)^\pm\big)}
\prod\limits_{j=1}^r\frac{c_j(l_j)^{r+1-j}\theta\big(c_j(l_j)^2\big)}
{c_j(k_j)^{r+1-j}\theta\big(c_j(k_j)^2\big)}
\\ \phantom{g_{\bf kl}=}{}
\times\prod\limits_{i=1}^r\frac{\prod\limits_{t=l_i+1}^{k_i}
\bigg\{\theta(c_i(t)b/c_1(k_1)\dotsm
 c_r(k_r))\prod\limits_{j=1}^r\theta(c_i(t)c_j(k_j))\bigg\}}
{\prod\limits_{t=l_i}^{k_i-1}
\bigg\{\theta(c_i(t)c_1(k_1)\dotsm c_r(k_r)/b)
\prod\limits_{j=1}^r\theta(c_i(t)/c_j(k_j))\bigg\}}
\end{gather*}
are mutually inverse.
\end{Theorem}



Note that the one-dimensional case of Theorem~\ref{cmit} is less
general than for Theorems \ref{amit} and \ref{bcmit}. Namely, it
corresponds to the special case when the sequences $a(t)$ and $c_1(t)$
are proportional. We have not been able to extend
Theorem~\ref{cmit} to the level of generality of the
other two inversions. Unfortunately, this means that our method of
proof of those results does not apply to Theorem~\ref{cmit}.
Instead we give another proof which
is more tedious in its details. The same method can be used to
give alternative proofs of
 Theorems~\ref{amit} and~\ref{bcmit}.

As before, to prove Theorem~\ref{cmit}
it is enough to verify \eqref{fg}
for ${\bf n}>\bf l=\bf 0$, which we state as a Lemma.

\begin{Lemma}\label{cfg}
Let $n_1,\dots,n_r$ be non-negative integers, not all equal to zero,
and let $b$ and
$c_j(k)$, $1\leq j\leq r$, $0\leq k\leq n_j$, be arbitrary scalars.
Then
\begin{gather*}
\sum_{k_1,\dots,k_r=0}^{n_1,\dots,n_r}
\prod\limits_{1\leq i<j\leq r}\frac{\theta(c_j(k_j)c_i(k_i))}{\theta(c_j(k_j)/c_i(k_i))}
\prod\limits_{j=1}^r\frac{\theta(c_j(k_j)b/c_1(k_1)\dotsm c_r(k_r))}{c_j(k_j)^{r+1-j}}
\\ \phantom{\sum_{k_1,\dots,k_r=0}^{n_1,\dots,n_r}}
\times\prod\limits_{i=1}^r\frac{\prod\limits_{t=1}^{n_i-1}
\bigg\{\theta(c_i(t)b/c_1(k_1)\dotsm
 c_r(k_r))\prod\limits_{j=1}^r\theta(c_i(t)c_j(k_j))\bigg\}}
{\prod\limits_{t=0,\,t\neq k_i}^{n_i}
\bigg\{\theta(c_i(t)c_1(k_1)\dotsm c_r(k_r)/b)
\prod\limits_{j=1}^r\theta(c_i(t)/c_j(k_j))\bigg\}}=0.
\end{gather*}
\end{Lemma}

\begin{proof}
 Pulling out all factors independent of $k_r$, our
sum can be written as
\begin{gather}\label{ss}
\sum_{k_1,\dots,k_{r-1}=0}^{n_1,\dots,n_{r-1}}A_{k_1,\dots,k_{r-1}}
\sum_{k_r=0}^{n_r}B_{k_1,\dots,k_{r}},
\end{gather}
where
\begin{gather*}
A_{k_1,\dots,k_{r-1}}
=\theta(b/c_1(k_1)\dotsm c_{r-1}(k_{r-1}))\prod\limits_{1\leq i<j\leq
 r-1}\frac{\theta(c_j(k_j)c_i(k_i))}{\theta(c_j(k_j)/c_i(k_i))}
 \\ \phantom{A_{k_1,\dots,k_{r-1}}=}{}
\times\prod\limits_{j=1}^{r-1}\frac{\prod\limits_{i=1}^r\prod\limits_{t=1}^{n_i-1}
\theta(c_i(t)c_j(k_j))}
{c_j(k_j)^{r+1-j}\prod\limits_{i=1}^{r-1}\prod\limits_{t=0,\,t\neq
 k_i}^{n_i}\theta(c_i(t)/c_j(k_j))\prod\limits_{t=0}^{n_r}\theta(c_r(t)/c_j(k_j))},
\\
B_{k_1,\dots,k_{r}}=
\frac{1}{c_r(k_r)}\prod\limits_{j=1}^{r-1}
\theta(c_j(k_j)b/c_1(k_1)\dotsm c_r(k_r),c_j(k_j)c_r(k_r))
\\ \phantom{B_{k_1,\dots,k_{r}}=}{}
\times\prod\limits_{i=1}^r\frac{\prod\limits_{t=1}^{n_i-1}\theta(c_i(t)b/c_1(k_1)\dotsm
 c_r(k_r),c_i(t)c_r(k_r))}{\prod\limits_{t=0,\,t\neq
 k_i}^{n_i}\theta(c_i(t)c_1(k_1)\dotsm
 c_r(k_r)/b,c_i(t)/c_r(k_r))}.
 \end{gather*}

We now rewrite the inner sum in \eqref{ss} using
 Gustafson's identity \eqref{cpf}. Replacing $a_k\mapsto
a_k/\sqrt\lambda$, $b_k\mapsto b_k\sqrt\lambda$ and using \eqref{ti},
\eqref{cpf} can be written as
\begin{gather}\label{cpfv}
\sum_{l=1}^k\frac{\prod\limits_{j=1}^{k-2}\theta(b_ja_l,b_j\lambda/a_l)}
{a_l\prod\limits_{j=1, j\neq l}^k\theta(a_ja_l/\lambda,a_j/a_l)}=0.
\end{gather}
Consider \eqref{cpfv} with $k=|{\bf n}|+1$ and parameters
\begin{gather*}
(a_1,\dots,a_{|{\bf n}|+1})=
\Big(c_1(0),\dots,\widehat{c_1(k_1)},\dots,c_1(n_1),\dots, c_{r-1}(0),\dots,\widehat{c_{r-1}(k_{r-1})},
\dots,
\\ \phantom{(a_1,\dots,a_{|{\bf n}|+1})=(}
c_{r-1}(n_{r-1}), c_r(0),\dots,c_r(n_r)\Big),
\end{gather*}
where the hats signify that the parameter is absent,
\begin{gather*}
(b_1,\dots,b_{|{\bf n}|-1})=
(c_1(1),\dots,c_1(n_1-1),\dots, c_r(1),\dots,c_r(n_r-1),
c_1(k_1),\dots,c_{r-1}(k_{r-1}))
\end{gather*}
and $\lambda=b/c_1(k_1)\dotsm c_{r-1}(k_{r-1})$. Then \eqref{cpfv}
splits naturally into $r$ parts, $\sum_{i=1}^rS_i=0$, where $S_1$ is
the sum of the first $n_1$ terms, $S_2$ the following $n_2$ terms and
so on until the final sum $S_r$ which has $n_r+1$ terms. It is
easy to check that $S_r$ equals the inner sum in \eqref{ss}, so that
\begin{gather*}
\sum_{k_r=0}^{n_r}B_{k_1,\dots,k_{r}}=-\sum_{i=1}^{r-1}S_i.
\end{gather*}

We will complete the proof by showing that
\begin{gather*}\sum_{k_1,\dots,k_{r-1}=0}^{n_1,\dots,n_{r-1}}A_{k_1,\dots,k_{r-1}}
S_i=0,\qquad i=1,\dots,r-1.\end{gather*}
 By a symmetry argument, it suffices to do this for $i=1$.
We have
\begin{gather*}S_1=\sum_{l=0,\,l\neq k_1}^{n_1}C_l, \end{gather*}
with
\begin{gather*}
C_l=
\frac{1}{c_1(l)}
\prod\limits_{j=1}^{r-1}\theta(c_j(k_j)b/c_1(l)c_1(k_1)\dotsm
 c_{r-1}(k_{r-1}),
 c_j(k_j)c_1(l))
 \\ \phantom{C_l=}{}
\times\frac{\prod\limits_{i=1}^r\prod\limits_{t=1}^{n_i-1}\theta
(c_i(t)b/c_1(l)c_1(k_1)\dotsm
 c_{r-1}(k_{r-1}),c_i(t)c_1(l))}
{\prod\limits_{i=1}^{r}\prod\limits_{t=0,\,t\notin \Lambda_i}^{n_i}
\theta(c_i(t)c_1(l) c_1(k_1)\dotsm
 c_{r-1}(k_{r-1})/b,c_i(t)/c_1(l))},
\end{gather*}
where
\begin{gather*}
\Lambda_i=
\begin{cases}\{k_1,l\}, & i=1,\\ \{k_i\}, &
 i=2,\dots,r-1,\\ \emptyset, & i=r.
\end{cases}
\end{gather*}

After some elementary manipulations, we obtain
\begin{gather*}
A_{k_1,\dots,k_{r-1}}C_{l}=
\frac{(-1)^r\theta(c_1(k_1)c_1(l),\gamma/c_1(k_1),
\gamma/c_1(l))}{c_1(k_1)^2c_1(l) \theta(c_1(l)/c_1(k_1))}
\prod\limits_{2\leq i<j\leq
 r-1}\frac{\theta(c_j(k_j)c_i(k_i))}{\theta(c_j(k_j)/c_i(k_i))}
 \\ \phantom{A_{k_1,\dots,k_{r-1}}C_{l}=}{}
\times\prod\limits_{j=2}^{r-1}\frac{\theta(c_j(k_j)c_1(k_1),c_j(k_j)c_1(l),c_j(k_j)\gamma/
c_1(k_1)c_1(l))\prod\limits_{i=1}^r\prod\limits_{t=1}^{n_i-1}
\theta(c_i(t)c_j(k_j))}{c_j(k_j)^{r+2-j}
\prod\limits_{t=0}^{n_1}\theta(c_1(t)/c_j(k_j))\prod\limits_{i=2}^{r}\prod\limits_{t=0,\,t\notin
 \Lambda_i}^{n_i}\theta(c_i(t)/c_j(k_j))}
 \\ \phantom{A_{k_1,\dots,k_{r-1}}C_{l}=}{}
\times\prod\limits_{i=1}^r\frac{\prod\limits_{t=1}^{n_i-1}
\theta(c_i(t)c_1(k_1),c_i(t)c_1(l),c_i(t)\gamma/c_1(k_1)c_1(l))}
{\prod\limits_{t=0,\,t\notin \Lambda_i}^{n_i}
\theta(c_i(t)/c_1(k_1),c_i(t)/c_1(l),c_i(t)c_1(k_1)c_1(l)/\gamma)},
\end{gather*}
where $\gamma=b/c_2(k_2)\dotsm c_r(k_r)$.
This is visibly anti-symmetric under interchanging
$k_1\leftrightarrow l$. Thus,
\begin{gather*}
\sum_{\substack{k_1,l=0\\k_1\neq l}}^{n_1}A_{k_1,\dots,k_{r-1}}C_{l}=0,
\end{gather*}
which completes the proof of the theorem.
\end{proof}
\begin{thebibliography}{99}
\bibitem[21]{gu}
Gustafson R.A., Multilateral summation theorems for ordinary and basic
 hypergeometric series in {${\rm U}(n)$}, \href{https://doi.org/10.1137/0518114}{\textit{SIAM~J. Math. Anal.}}
 \textbf{18} (1987), 1576--1596.

\bibitem[40]{r}
Rosengren H., Elliptic hypergeometric series on root systems, \href{https://doi.org/10.1016/S0001-8708(03)00071-9}{\textit{Adv.
 Math.}} \textbf{181} (2004), 417--447, \href{https://arxiv.org/abs/math.CA/0207046}{arXiv:math.CA/0207046}.

\bibitem[46]{s}
Schlosser M.J., Multidimensional matrix inversions and {$A_r$} and {$D_r$}
 basic hypergeometric series, \href{https://doi.org/10.1023/A:1009705129155}{\textit{Rama\-nujan~J.}} \textbf{1} (1997),
 243--274.

\end{thebibliography}
\end{document}
